% =============================================================
% Section: Policy implications
% =============================================================

\section{Implications for reform design}\label{sec:conclusion}

The relevant policy object is no longer only the original 36-hour version of PEC 8/2025. The live debate also includes 40-hour, 5x2, and phased alternatives. The model therefore reports the full hours-cap schedule. A 40-hour cap has a much smaller output-neutrality benchmark than the endpoint and leaves the representative agent GHH diagnostic positive. A single-step move from 44 to 36 hours is different. It affects most Brazilian formal workers and reduces the labor input of the economy by about $14\%$ mechanically. At that endpoint, the output loss is offset by a one-off productivity gain of $6.63\%$ under one-sided fatigue and $8.18\%$ under the two-sided penalty.

The 36-hour endpoint is large relative to Brazil's recent productivity record. The PWT comparison should be read as a scale benchmark, because an annual macro residual is not the same object as the short-run productivity response in the model. It is nevertheless informative about magnitudes: post-1990 Brazilian TFP growth has been weak, so the endpoint would require an unusual acceleration if output were to be preserved over a short implementation horizon. Without such a productivity response, the model implies a short-run output loss of $6.6$ to $8.0\%$ and a decline in the aggregate GHH diagnostic of $1.8$ to $3.6\%$ at the endpoint.

A moderate reduction can raise the model's GHH diagnostic. In the representative agent calibration, cutting the workweek from $44$ to $41$ or $42$ hours raises that indicator because the value assigned to extra leisure exceeds the consumption loss. The 40-hour alternative is near the edge of the positive region, and below $40$h the calculation turns negative. This result depends on the assumed hours-productivity schedule and representative preferences. It is not a social welfare ordering over heterogeneous workers.

The cost of the endpoint is concentrated. Small firms bear the largest per-worker productivity requirement. Industry and agriculture carry the largest per-sector requirements, and informality rises in all three sectors. Services account for most of the aggregate output loss through their employment and output share, not through a larger per-worker effect. The sectoral exercise is partial-equilibrium and sector-by-sector rather than an input-output incidence model, but it is enough to show that the adjustment burden is uneven.

Informality rises but does not dominate the output cost. The calibrated model, disciplined by the observed formal-informal wage premium, finds an aggregate informality increase of $1.6$ to $1.9$ pp. That rise matters for welfare accounting, especially for workers pushed into informality. It is small relative to the mechanical hours contraction. The output cost comes primarily from fewer hours worked.

These conclusions have limits. The model is short-run and partial-equilibrium. Endogenous capital adjustment over a three-to-five-year horizon lowers the required productivity gain by about one to two percentage points. Total employment is fixed, which shuts down work-sharing. Evidence on work-sharing from Germany \citep{Hunt1996} and Brazil's 1988 constitutional reform \citep{GonzagaMenezesFilhoCamargo2003} suggests that this restriction is compatible with available evidence, but it is still a restriction. Compliance, overtime, organizational disruption, wage bargaining, and informality responses could move the relevant benchmark in either direction. The Portuguese 1996 experience \citep{AsaiLopesTondini2024} documents firm-level productivity gains larger than the model produces, which suggests that the model is conservative on the hourly-productivity response.

Under these limits, the productivity arithmetic suggests that phased paths deserve separate evaluation. An intermediate stop near $40$ to $42$ hours keeps the productivity requirement far below the 36-hour endpoint and gives capital, informality, and firm-level productivity time to adjust before a further step. The implication is not that the workweek should remain at 44 hours, nor that the model identifies the socially optimal cap. It is that the size and timing of the reduction matter for the tradeoff between output neutrality, representative agent welfare, and incidence across firms and workers.
